\chapter{The informed investor}\label{ch:trading}

\begin{chapterplan}
What does an informed order earn when the price already counts the entry that it brings? Which two features of
the auction matter, and how do they move with the incumbent's strength?
\end{chapterplan}

\section{The gross return of an order}\label{sec:trading-returns}

If the investor places the order \(q\), it trades at \(P(q)\) and, given its class \(j\), expects
\begin{equation}\label{eq:trading-objective}
U_j(q)=q\,\bigl[t_0+e(q)\,(t_j-t_0)-P(q)\bigr]-k|q| .
\end{equation}
The price is the market makers' average of the same payoff over the two classes (Chapter~\ref{ch:beliefs}):
\(P(x)=t_0+e(x)\,[t_L-t_0+\DT\,\mu(x)]\).

\begin{paperresult}{Lemma A.3 (Informed trading)}
Let \(\DT\ge0\). Fix a tuple that satisfies (E1), (E4) and (E6), with entry \(e(x)\in[\rho,1]\) after every order.
Conditional on its arrival, the investor in each class earns
\[
\begin{aligned}
U_H(1)&=e(1)\DT[1-\mu(1)]-k,& U_H(0)&=0,& U_H(-1)&=-e(-1)\DT[1-\mu(-1)]-k,\\
U_L(-1)&=e(-1)\DT\,\mu(-1)-k,& U_L(0)&=0,& U_L(1)&=-e(1)\DT\,\mu(1)-k.
\end{aligned}
\]
Write \(K_H=e(1)\DT[1-\mu(1)]\) and \(K_L=e(-1)\DT\,\mu(-1)\) for the gross returns of the two correct orders.
\begin{enumerate}[label=(\alph*)]
\item A wrong-signed order earns at most \(-k\), so abstention strictly beats it.
\item Class \(H\) buys if \(K_H>k\) and abstains if \(K_H<k\). At \(K_H=k\), the purchase and abstention are its two
best orders. Class \(L\) sells if \(K_L>k\) and abstains if \(K_L<k\). At \(K_L=k\), the sale and abstention are
its two best orders.
\item Under every candidate profile, \(\rho\DT[1-M(\pi)]\le K_H\le\DT\) and \(\rho\DT\,m(\pi)\le K_L\le\DT\).
Hence \(\rho\,\kappa(\pi)\DT\le K_j\le\DT\), where \(\kappa(\pi)=\min\{m(\pi),1-M(\pi)\}\).
\item \(K_H\) depends on the candidate only through \(\mu(1)\) and \(e(1)\), and \(K_L\) only through \(\mu(-1)\)
and \(e(-1)\).
\end{enumerate}
\end{paperresult}

\emph{Proof of the payoffs.} Substitute the price into the bracket of \eqref{eq:trading-objective}:
\begin{align*}
t_0+e(q)(t_H-t_0)-P(q)&=e(q)\bigl[(t_H-t_0)-(t_L-t_0)-\DT\mu(q)\bigr]=e(q)\,\DT\,[1-\mu(q)],\\
t_0+e(q)(t_L-t_0)-P(q)&=-e(q)\,\DT\,\mu(q). \why{\(t_0\) and \(t_L-t_0\) cancel}
\end{align*}
Multiply by \(q\) and subtract \(k|q|\); this gives the six payoffs. In words: the gross return is entry times the spread times the residual uncertainty about the class. The investor
receives \(t_0\) and the low-class gain at the end and pays them in the price, so they net out.

\paragraph{Best orders.}\label{sec:trading-best}\label{sec:trading-bound}
\emph{(a)} The gross part of a wrong-signed order, \(-e\DT[1-\mu]\) or \(-e\DT\mu\), is never positive, so its
payoff is at most \(-k<0\). \emph{(b)} So each class compares its correct order with abstention:
\(U_H(1)-U_H(0)=K_H-k\) and \(U_L(-1)-U_L(0)=K_L-k\). At equality the tie rule decides. \emph{(d)} is visible in \(K_H\) and \(K_L\). Abstention earns zero whatever happens after the zero order. So the problem splits into
a purchase side for class \(H\) and a sale side for class \(L\); Chapter~\ref{ch:coexistence} solves them one at a
time.

\paragraph{The trading bound.} \emph{Proof of (c).} Entry lies in \([\rho,1]\), and the posterior after a nonzero order lies in
\([m(\pi),M(\pi)]\) for every \(0<\beta<1/2\) (Lemma A.1(b)). All factors are nonnegative, so
\begin{align*}
K_H&=e(1)\DT[1-\mu(1)]\ge\rho\DT[1-M(\pi)]\ge\rho\,\kappa(\pi)\DT, &
K_L&=e(-1)\DT\mu(-1)\ge\rho\DT\,m(\pi)\ge\rho\,\kappa(\pi)\DT .
\end{align*}
The ceiling uses \(e\le1\) and \(0\le\mu\le1\). \qed
Since \(1-M(\pi)=m(1-\pi)\) and \(m\) rises strictly, \(\kappa(\pi)=m(\min\{\pi,1-\pi\})\le1/(2+u)\), with the
peak at \(\pi=\tfrac12\). The class that the prior favors binds, because the price already leans toward it.
The bound holds in every candidate; its floor is positive only because \(\rho>0\). If \(\DT<k\), every class abstains in every equilibrium. If \(\rho\,\kappa(\pi)\DT>k\), class \(H\) buys and class
\(L\) sells for sure. Chapter~\ref{ch:coexistence} studies the region between.

\section{The two numbers}\label{sec:trading-twonumbers}

\begin{paperresult}{Lemma A.4 (The auction through the spread and the threshold)}
Fix the model's independence and information structure, a prior \(\pi\in(0,1)\), the two-point preparation cost
with \(\rho\in(0,1)\) and \(c_H>0\), a trading cost \(k>0\), and one pair of tie rules: an indifferent investor may
mix over its best orders or abstains, and an indifferent high-cost challenger prepares or may prepare with any
probability. At fixed presence, fix \(0<\lambda<1\); with research, presence is \(\hat\lambda\). Fix
\(0<\beta<1/2\). Consider two cash second-price auctions with finite-mean classes satisfying the class ranking,
continuous incumbent laws, \(\DT>0\) and \(\gamma=g_H-g_L>0\), and the same spread and preparation threshold
\(\tau=(c_H-g_L)/\gamma\). Their equilibrium sets projected onto \((\sigma_H,\sigma_L,\mu,e)\) coincide.

The common conditions are Bayes' rule, the threshold entry rule under the chosen tie rule, and optimal orders at
\[
\begin{aligned}
U_H(1)&=e(1)\DT[1-\mu(1)]-k,& U_H(-1)&=-e(-1)\DT[1-\mu(-1)]-k,\\
U_L(-1)&=e(-1)\DT\mu(-1)-k,& U_L(1)&=-e(1)\DT\mu(1)-k,
\end{aligned}
\qquad U_j(0)=0.
\]
Each auction gives prices and price beliefs by
\(P(x)=t_0+e(x)[w_L+\DT\mu(x)]\), \(\zeta(P(x))=[e(x)-\rho]/(1-\rho)\), \(\bar\mu(P(x))=\mu(x)\).
With research, fix also the quadratic cost coefficient \(a>0\) and cap \(0<\lb<1\), and require optimal orders
after a hypothetical success even at zero research. The equilibrium sets projected onto
\((\hat\lambda,\sigma_H,\sigma_L,\mu,e)\) coincide. The common conditions also include
\(\hat\lambda=\min\{\lb,\Psi/a\}\), \(\Psi=\pi\max_q U_H(q)+(1-\pi)\max_q U_L(q)\), with Bayes' rule evaluated at \(\hat\lambda\). These conclusions hold under all four tie rules.
\end{paperresult}

\emph{Proof.} \emph{Forward.} Lemma A.2(c) makes the belief at \(P(x)\) equal
\(\mu(x)\). Then \(B(\mu)-c_H=\gamma(\mu-\tau)\), so (E3) is the threshold rule in \(\tau\), \(\rho\) and the tie
rule. Lemma A.3 gives payoffs in \(\DT\) alone. \emph{Back.} Threshold
entry does not fall with the posterior, so the claim on constructed tuples (Chapter~\ref{ch:beliefs}) builds each
auction's \(P\), \(\zeta\) and \(\bar\mu\), and (E1) to (E6) hold by substitution. \emph{Research.} The payoff
\(\tilde\lambda\Psi-(a/2)\tilde\lambda^2\) is strictly concave, so \(\hat\lambda=\min\{\lb,\Psi/a\}\). \qed

\section{The two numbers along a path}\label{sec:trading-path}

In words: a stronger incumbent moves the class value gap into target proceeds, so the spread and the threshold
race. \(J\), the return of a purchase that just recruits, can go either way.

\begin{paperresult}{Lemma A.5 (The two numbers along a path)}
Fix ranked finite-mean classes, the reserve and the two-point preparation cost. A path is a family \(F_r\) of
continuous incumbent laws indexed by strength on an interval, with \(F_{r'}\le F_r\) when \(r'>r\). Write
\(\DT(r)=\DT(F_r)\), \(\gamma(r)=g_H(F_r)-g_L(F_r)\), and, where \(\gamma(r)>0\),
\(\tau(r)=[c_H-g_L(F_r)]/\gamma(r)\). The class mean gap is fixed:
\(\bar\Gamma=\E_H\theta-\E_L\theta=\DT(r)+\gamma(r)\).
Each \(g_j\) and \(\gamma\) is nonincreasing, and \(\DT\) is nondecreasing. For \(r'>r\),
\[
\begin{aligned}
\DT(r')-\DT(r)&=\int_p^\infty[F_r(y)-F_{r'}(y)]D(y)\,dy,\\
g_j(F_r)-g_j(F_{r'})&=\int_p^\infty[F_r(y)-F_{r'}(y)][1-G_j(y)]\,dy.
\end{aligned}
\]
Each change is strict exactly when its integral is positive. Where \(\gamma>0\), \(\DT>0\) and \(\tau\ge0\), both
\(\tau\) and \(\DT\tau\) are nondecreasing. Their exact increments are
\[
\tau(r')-\tau(r)=\frac{g_L(F_r)-g_L(F_{r'})+\tau(r)[\DT(r')-\DT(r)]}{\gamma(r')},
\]
and \(\DT(r')\tau(r')-\DT(r)\tau(r)=\DT(r)[\tau(r')-\tau(r)]+\tau(r')[\DT(r')-\DT(r)]\).
Thus the threshold rises strictly exactly when the numerator in its increment is positive; the product rises
strictly exactly when the right side of its increment is positive.

On a differentiable path the return \(J=\DT(1-\tau)\) obeys
\(J'=[\bar\Gamma(g_H-c_H)\DT'+\DT\gamma g_H']/\gamma^2\).
Its direction depends on the balance between spread growth and the decline in high-class profit. The paper gives an
instance with point classes \(h\) and \(\ell\), a reserve \(p\), a high cost \(c_H\), and \(F_r\) uniform on a short
interval centered at \(r\). The strengths \(r\) run over an interval on which every incumbent value lies strictly
between \(\ell\) and \(h\). There
\(\DT=r-\ell\), \(g_L=0\), \(g_H=h-r\), \(\tau=c_H/(h-r)\),
and the return \(J\) is lower at the right end of the interval of strengths than at the left end. In the
uniform-incumbent example, \(J\) rises on \([r_0,r_2]\) by the uniform case of Proposition~3(c).
\end{paperresult}

\emph{Proof.} Every value meets the reserve, so \(t_j+g_j=\E_j\theta\) and \(\DT+\gamma=\bar\Gamma\). The profit integral \(g_j(F)=\int_p^\infty F(y)[1-G_j(y)]\,dy\) gives the change in \(g_j\), and
\(\gamma=\int_p^\infty F D\,dy\) gives \(\DT(r')-\DT(r)=\gamma(r)-\gamma(r')\ge0\). For the threshold, use
\(\gamma(r')=\gamma(r)-[\DT(r')-\DT(r)]\) and \(\tau(r)\gamma(r)=c_H-g_L(F_r)\):
\begin{align*}
\tau(r')-\tau(r)=\frac{c_H-g_L(F_{r'})-\tau(r)\gamma(r')}{\gamma(r')}
=\frac{g_L(F_r)-g_L(F_{r'})+\tau(r)[\DT(r')-\DT(r)]}{\gamma(r')}\ge0 .
\end{align*}
For \(J=\DT(g_H-c_H)/\gamma\), differentiate with \(\gamma'=-\DT'\) and \(\gamma+\DT=\bar\Gamma\). \qed

\begin{watchout}
The threshold result needs \(\gamma>0\) and \(g_L\le c_H\), that is \(\tau\ge0\). Without it the threshold can fall.
\end{watchout}


\begin{chaptersummary}
The investor earns only on the class: \(K_H=e(1)\DT[1-\mu(1)]\), \(K_L=e(-1)\DT\,\mu(-1)\). The trading bound \(\rho\,\kappa(\pi)\DT\le K_H,K_L\le\DT\) holds in every candidate.
Orders, posteriors, entry and research see the auction only through \(\DT\) and \(\tau\). A stronger incumbent
never narrows the spread and, while \(g_L\le c_H\), never lowers the threshold.
\end{chaptersummary}
